\section{Introduction}
\label{chapter:01}

% block-id: 01.introduction.heading
\hypertarget{01.introduction.heading}{}
\subsection{High-frequency loss and current sharing}

% block-id: 01.introduction.p1
\hypertarget{01.introduction.p1}{}
Litz wire controls high-frequency current through strand insulation and longitudinal transposition. Strand diameter sets the local skin-effect scale, neighboring conductors and the return shape proximity fields, and common terminals couple the complex strand currents. A structural comparison therefore requires both an electromagnetic model and consistent definitions of copper area, axial length, and terminal conditions.

% block-id: 01.introduction.p2
\hypertarget{01.introduction.p2}{}
Discrete-transposition modeling with two-dimensional FEM (DTM) \cite{P01} connects strand-position permutations to cross-sectional fields. Its extension to nonideal compound transposition (iDTM) \cite{N01} treats hierarchical and subbundle effects. Improved segmented equivalent circuits \cite{P02}, FEM–PEEC coupling \cite{P03}, and quasi-two-dimensional multiconductor formulations \cite{P04} use different reductions of field and circuit behavior. Continuum models \cite{P05} and complex-permeability homogenization \cite{P09} reduce the cost of winding calculations. Their results must be interpreted through the fields, current constraints, and boundary conditions retained by each method.

% block-id: 01.introduction.p3
\hypertarget{01.introduction.p3}{}
This work connects three levels through reviewable numerical verification: round-wire analysis, multistrand two-dimensional common-terminal coupling, and finite-length longitudinal PEEC. The analytical solution supplies an internal-impedance reference; the two-dimensional model solves both intrastand fields and complex strand currents; the longitudinal model examines the influence of the finite-length kernel and integration settings on resistance. Traditional multilevel Litz provides the baseline for subsequent structural comparisons, while cross-layer paths are candidate geometries to be evaluated. The numerical conclusions apply to the degrees of freedom and boundary conditions retained at each level.

% block-id: 01.scope.heading
\hypertarget{01.scope.heading}{}
\subsection{Research questions and contributions}

% block-id: 01.scope.p1
\hypertarget{01.scope.p1}{}
This paper addresses three related questions: whether independent discretizations reproduce round-wire skin current and loss; how uniform complex strand currents can be characterized under common terminal voltage; and how finite-length longitudinal PEEC resistance responds to length, integration, and cross-sectional discretization. We organize verification by observable, derive a current-sharing condition for a specified linear field–circuit model, and report quantitative results for a 640-strand two-dimensional case and longitudinal cases with 1264 and 5024 degrees of freedom. Supplementary three-dimensional imposed-field, potential-diagnostic, and interface cases identify the implementation's applicability and the conditions required for further research.
